% Recuperación de 02_trit_geometrias y de la holonomía con memoria.
\subsubsection{Paramétrisation temporelle de la compatibilité et orientation tritique}
\label{trit-temporal-compatibilidad}

L’interprétation temporelle du TRIT est établie sur les transitions de
l’état complet. Soit \(x\xrightarrow{\rm adm}y\) la relation d’admissibilité déterminée
par les opérations du TPK dans \(\widehat X\). Une trajectoire paramétrée
par un ensemble discrètement ordonné \(\mathcal T\) est une application
\[
 \gamma:\mathcal T\longrightarrow\widehat X,\qquad
 \gamma(t)\xrightarrow{\rm adm}\gamma(t^+),
\]
où \(t^+\) est le successeur de \(t\). La relation de compatibilité précède
le choix du paramètre : l’ordre permet d’exprimer ses transitions comme
une évolution. Lorsque le terme section est utilisé, on déclare en outre la
projection \(p:E\to\mathcal T\) et la condition
\(p\circ\gamma=\operatorname{id}_{\mathcal T}\). Sont ainsi spécifiés
l’espace des états, l’ordre temporel et l’information transportée.

La signature \(\tau\) sélectionne un régime quadratique et admet une lecture
temporelle relative à cette paramétrisation. Le type \(+1\), avec
\(J_{+1}^{2}=-I\), représente le retour et la mémoire ; le type \(0\), avec
\(J_0^{2}=0\), représente le seuil de transition ; le type \(-1\), avec
\(J_{-1}^{2}=I\), représente l’ouverture dans deux directions propres. Dans la
terminologie temporelle du modèle, ces fonctions correspondent,
respectivement, au passé, au présent et au futur. La correspondance concerne
la position opératoire d’une transition dans une histoire :
antécédent conservé, actualisation et prolongement admissible.

Cette interprétation acquiert un contenu par trois opérations distinctes.
La restriction d’une trajectoire restitue le préfixe déjà construit.
L’actualisation \(\mathcal U_t\) compose sélection, transport et enregistrement dans
la transition observée. Le prolongement considère les continuations qui
conservent les conditions de frontière de ce préfixe. Si
\(\mathcal P_n(x)\) désigne l’ensemble des continuations admissibles de
longueur \(n\) depuis \(x\), la suppression du dernier pas définit
\[
 r_n:\mathcal P_{n+1}(x)\longrightarrow\mathcal P_n(x).
\]
Une continuation illimitée est une collection \((\gamma_n)_{n\geq1}\) qui
satisfait \(r_n(\gamma_{n+1})=\gamma_n\). Le système projectif conserve
ainsi les restrictions que chaque profondeur impose aux
précédentes. L'existence d'une collection compatible est établie dans le
domaine de prolongement correspondant ; la seule définition des
ensembles \(\mathcal P_n(x)\) ne remplace pas cette démonstration.

La bidirectionnalité combine donc la reconstruction des antécédents
avec la compatibilité des continuations. Sur la suite des signatures,
l’involution \(C_{\rm rev}\) inverse l’ordre et change le signe de chaque
composante. Sur la trajectoire agissent également les transformations de
position, de feuillet et de frontière spécifiées par le transport. La signature
inversée est une observation de la trajectoire conjuguée, non un remplacement
de ses autres coordonnées. Le seuil demeure distingué : la valeur
visible \(0\) conserve sa fonction de transition même lorsque son quotient
ou son registre historique sont distincts de zéro.

La relation avec les géométries exponentielles est tout aussi précise. Dans
chaque régime fixé, \(\exp(tJ_\tau)\) compose les paramètres et conserve la
norme algébrique. La clôture elliptique appartient à cette représentation locale.
Le retour de phase de \(\Gamma_9\), en revanche, agit sur la trajectoire
avec mémoire, conformément à \eqref{eq:holonomia-memoria}. Une trajectoire peut
retrouver sa phase et conserver un nouvel antécédent : la donnée historique a
changé bien que l’observation de phase soit la même. Le seuil parabolique
retient un déplacement linéaire ; l’ouverture hyperbolique sépare deux
directions propres d’expansion et de contraction. Ces trois réalisations
fournissent des lois de composition, tandis que le TPK détermine leur
sélection et leur concaténation effective.

La suspension apériodique de la section \ref{sec:generacion} explicitera
cette distinction au moyen d’un facteur de phase fini, d’une rotation irrationnelle
et d’un degré de mémoire. Le lien entre TRIT et temporalité s’exprime ainsi
par des opérations sur les trajectoires et par leurs représentations
quadratiques. Son contenu dépasse une classification des signes : il conserve le
régime, le sens du transport et la relation entre antécédent,
transition et continuation.
